\section{Introduction}
\label{sec:v15-introduction}
Comparison of statistical experiments becomes a different problem when
the comparison channel must itself run with a prescribed finite memory.
A causal channel may preserve the filtration and nevertheless require
more persistent state than the experiment's consumer possesses. Moreover,
randomizing among complete designs can enlarge the achievable set even
when each design uses the same number of labels. The information exposed
by the random selector, and the cost of retaining its value, are therefore
part of the mathematical object being compared.

We study marked causal experiments with this resource structure. A mark
is an actual random variable coupled to the source, not an independent
copy with the same marginal law. The intrinsic deficiency is the smallest
worst-feedback variation error among simulators at the declared state
and randomization budget. Private, hidden-selector and visible-selector
budgets lead to distinct feasible sets. The private image is generally
nonconvex. Its exact finite certificates must consequently test that
image itself rather than replace it by the convex hull of deterministic
implementations.

The main result is a same-budget physical certification theorem. It
starts with an active microscopic acquisition, passes through a finite
trial presentation, quantizes its reports, constructs a coherent rational
marked experiment, and transfers finite private certificates back to
the original physical experiment. Each presentation map is stateless.
The register cost paid by the final physical simulator is therefore the
original finite budget, without an unpriced model-selection index.
The analytic error and the finite optimization error are both explicit.

\subsection{The principal theorem}
Let $b$ be a fixed private or finite-selector signature. For a pair of
finite-horizon active Gaussian experiments $E,F$, write
$\delta_b(E,F)$ for their intrinsic marked feedback deficiency. Let
$E_{\mathbb Q},F_{\mathbb Q}$ be finite rational marked presentations
constructed by the procedures of this paper. Suppose their two-sided
presentation errors are at most $a,c$, respectively. If finite local
testing and an executable same-budget row table give bounds $L,U$ for
$\delta_b(E_{\mathbb Q},F_{\mathbb Q})$, then
\begin{equation}\label{eq:v15-intro}
 \max\{0,L-a-c\}\le\delta_b(E,F)\le\min\{1,U+a+c\}.
\end{equation}
Theorem~\ref{thm:v15-main} proves this statement and a convergent
construction under effective graph and marked-integration data. The
lower bound remains valid without assuming attainment for the original
nonfinite experiment. The upper bound comes with an actual physical
simulator. The resulting intervals can be nested and made arbitrarily
short. This is a finite-horizon effective theorem with specified analytic
inputs, not a computability assertion for arbitrary measurable models.

There are three substantive steps. First, report truncation and histogram
reconstruction are compared before controller optimization, using the
original preparation and a coupled actual mark. This avoids unjustified
bounds on feedback-conditioned densities. Second, rationalization is
performed on joint marked prefix masses. A weighted normalization
inequality cancels the small-prefix denominator. It gives a coherent
rational causal tree and a two-sided feedback error without deciding
which original prefix probabilities vanish. Third, piecewise affine
graph-core paths give enriched residual bounds whose values themselves
converge to zero. The intervention is applied after this domain-controlled
path calculation; no invariance of the generator domain under kicks
is assumed.

Two further results connect the finite geometry to state complexity.
Theorem~\ref{thm:v15-witness} bounds the number of distinct testing
witnesses by the affine dimension of the attainable behavior image and
the strict error gap. It keeps the original nonconvex private image;
its bound is on tests, not the number of policy cells or computational
running time. Theorem~\ref{thm:v15-task} gives packing lower bounds and
covering upper bounds for the positive-error width of a finite family
of marked prediction tasks at a source-consuming restart. Completing
that family recovers the universal marked continuation state at zero
error. Causal update constraints remain necessary for simultaneous
checkpoint realizations.

The effective data are not left as an integration oracle.
Theorem~\ref{thm:v15-constructive-core} constructs an enumerable
microscopic graph core and verifies its Gram and marked probability
operations by certified full-measure collision-chart exhaustion. It
applies to fixed-particle hard spheres with algebraic velocity kicks
and bounded computable input data, without requiring conserved
observables. No particle-number-uniform rate is inferred.

The physical data also admit explicit closed-form instances.
Theorem~\ref{thm:v15-quarter} gives an exact deficiency $1/4$ for a
continuous Gaussian acquisition of a conserved momentum coordinate in
an actively intervened colliding hard-sphere system. Its trial Gram
data are rational and exact. Proposition~\ref{prop:v15-bump} then gives
a collision-compatible localized position observable with a genuinely
nonzero graph residual $R=192$ and a rigorously bounded physical error.
The conserved-sector family admits the complete finite-rational
procedure at every fixed horizon, including rigorous scalar integration.

\subsection{Classical antecedents and the distinction of resources}
The general comparison of experiments and its risk interpretation are
classical \cite{Blackwell1953,LeCam1964}. Filtered randomization criteria
predate the present finite-state formulation
\cite{Norberg2002,Weisshaupt2002}. In particular, Weisshaupt's Lemma~1
and Theorem~5 use a compact convex set of filtered stochastic operators
with values in finitely additive measures; their proof reduces the
filtered statement to convex separation. They do not price a fixed
persistent register. That comparison cannot be applied to our generally
nonconvex private image by silently convexifying it. We retain the
classical result as an antecedent, not a theorem displaced by this paper.

Common-information control and real-time coding provide essential
information-state precedents
\cite{NayyarMahajanTeneketzis2013,Witsenhausen1979}. Compatible-cover
language has an antecedent in incomplete-machine minimization
\cite{PaullUnger1959}. The distinction here is the stated marked,
resource-constrained comparison problem and its error transfer, not the
invention of dynamic programming or compatible covers. A complete
proof-level priority comparison with the original Norberg and
Paull--Unger texts remains to be established; we make no assertion of
exhaustive priority over those works.

Separate-affine vertex interpolation, subdivision, normed-space packing,
real quantifier elimination \cite{Basu2014}, variation of constants,
and graph-core approximation \cite{Trotter1958} are also classical.
Their roles in the proofs are identified where used. The contribution
claimed here is the original-budget physical certificate chain,
including its actual marked feedback bridges, zero-prefix-stable
rationalization and convergent enriched certificates, together with
the stated resource-image and task-resolution consequences. None of
these claims relies on counting source files or passing a diagnostic
suite as a mathematical argument.

\subsection{Organization and the program of applications}
Section~\ref{sec:v13-deficiency} develops intrinsic resource comparison,
its finite local certificates and the new behavior-dimension witnesses.
Section~\ref{sec:v13-controlled} treats endogenous common decisions and
positive-error restart complexity. Section~\ref{sec:v13-regeneration}
retains the stationary regenerative theorem with its actual physical
and register reset. The active microscopic construction and the original
Galerkin residual inequalities follow. Sections~\ref{sec:v15-finite}
and \ref{sec:v15-enriched} supply the two missing approximation steps;
Section~\ref{sec:v15-effective-physical} constructs effective
microscopic data; Section~\ref{sec:v15-certification} proves the joint theorem and
Section~\ref{sec:v15-physical-data} instantiates certified physical data.

The associated eleven-component program is a dependency graph. Its
primary geometric A2 branch is independent of the present comparison
theory; an acquisition adapter is a different relation from theorem
dependence. The present microscopic B4 edge proves fixed-horizon marked
certification. It does not turn $L^2$ transport estimates into the
particle-number-uniform nonlinear action-sublevel corrector of the
historical kinetic program. Those larger targets and their derivations
are preserved in the source history and accompanying complete development.
This article states precisely the new implications, without using
repository order as a substitute for proof.
